% Section 5 -- Results (outline: paragraphs and what each must say; no prose yet)
% Structure (28 October): an opening and two subsections, one per explanation
% of Section 2: the payoff horizon (5.1), then the window (5.2).
%
% Numbers from the draft exhibits of the replication package (R0, n = 65,536,
% CPU unless noted); none of the tables is in the paper yet:
%   lifecycle     output/figures/fig-lifecycle-fit.csv        (fig:lifecycle-fit)
%   horizon       output/tables/tab-payoff-horizon.csv        (tab:payoff-horizon)
%   flows and     output/tables/tab-flows-accounting.csv      (tab:flows-accounting;
%   accounting                                                 merged 28 October)
%   mechanisms    output/tables/tab-mechanisms.csv            (tab:mechanisms; GPU
%                 cells of the 30 September batch, decay rows on CPU)
%   varrho 0.96   data/model/r0/lifecycle/*-varrho0.96.csv    (flows, horizon)
% The search box is in data/model/r0/search-box.json.

\outpara{Opening: the question and the answer (before the first subsection)}
\begin{itemize}
  \item The question from Section~\ref{sec:data}: do older owners live in
    less efficient homes because efficiency pays off over fewer years
    (the payoff horizon), or because they buy homes less often (the
    window)? Section~\ref{sec:model} gave the direction of each force;
    this section gives their size at the estimates.
  \item The answer up front: almost all of the gradient comes from
    purchases. The payoff horizon is short at every age, because
    efficiency wears off, so it hardly varies with age. Older owners buy
    less often, and buying an efficient home is the main way to get one.
  \item Road map: the horizon (5.1), then the window (5.2).
  \item \open{Title: ``Results'' or one that states the question, e.g.
    ``Why Older Owners Live in Less Efficient Homes''.}
\end{itemize}

\outpara{Opening: the life cycle at the estimates (Figure fig:lifecycle-fit)}
\begin{itemize}
  \item The model's profiles of homeownership, the A--C share, mortgages,
    income, and consumption by age. Decided 12 and 21 October: model only,
    opens Section 5.
  \item What the reader needs from it: when households buy (renters buy in
    their late twenties and thirties, few purchases after 50), and the
    A--C share falling with age.
  \item Outcomes, not trade-offs: this is where sentences such as ``young
    households rent and save, then buy'' belong (kept out of
    Section~\ref{sec:model} on purpose).
\end{itemize}

\subsection{The payoff horizon}

\outpara{The horizon by age (Table tab:payoff-horizon)}
\begin{itemize}
  \item The horizon of equation~\eqref{eq:payoff-horizon}, built up one
    force at a time, at ages 25, 45, 65, 75.
  \item Survival alone: 26.9, 22.0, 14.6, 10.0 years. This is the intuition
    of the first explanation, and it falls steeply with age.
  \item Adding loss of efficiency: 5.3, 5.3, 5.0, 4.4. Adding moving at the
    model's rate: 3.4, 3.0, 3.1, 2.9. Loss of efficiency caps the horizon at
    every age, so survival hardly matters before 75.
  \item With the moving rates of owners in UKHLS the horizon rises in
    midlife (3.8, 4.7, 4.6, 4.2): the old move less, which offsets the
    higher probability of dying (Section~\ref{sec:model}).
\end{itemize}

\outpara{What a flat horizon implies}
\begin{itemize}
  \item Proposition~\ref{prop:window}: the gap between the two thresholds is
    the extra cost of renovating in place spread over $\widetilde H$. With
    $\widetilde H$ nearly flat in age, the gap is about as wide at 75 as at
    35. Proposition~\ref{prop:buy}: a buyer's choice of an efficient home
    hardly depends on age.
  \item So the margins on which the horizon acts (what a buyer chooses,
    whether an owner renovates in place) should be flat in age, and the
    horizon cannot produce much of a gradient. Section 5.2 checks this in
    the accounting.
  \item Optional: the reference-home example parked from Section 3.3 (bill
    saving against the cost of renovating) fits here, now that the
    estimates are known (parked-text.md).
\end{itemize}

\outpara{The payoff per year and how long it lasts (Table tab:mechanisms,
Panel A)}
\begin{itemize}
  \item Change one ingredient at a time, let households re-optimize, keep
    the other parameters at the estimates. Report the A--C share by age and
    the gap and the ratio between 25--34 and 65--75.
  \item Comfort (the payoff per year): none or doubled shifts every share
    (gap 7.4 and 11.3) but hardly the ratio (1.76 and 1.68). A larger
    payoff per year raises the level, not the age profile.
  \item Retention (how long the payoff lasts): $\varrho=0.88$, 0.92, 0.96
    raises every share; the gap widens and then narrows (11.4, 11.5, 6.2)
    and the ratio falls (1.68, 1.49, 1.17). A longer horizon does not
    steepen the gradient: with slow decay the stock remembers older
    purchases, so older owners catch up.
  \item \open{$\eta=0$ is the lower edge of the training box [0, 0.15],
    the same concern as for $\varrho$ at 0.80 and 0.99. Keep only the
    doubled cell (interior), add an interior cell such as $\eta/2$ (new
    run), or say so in the note.}
\end{itemize}

\outpara{When efficiency lasts 25 years: $\varrho=0.96$ (promised in
Section~\ref{sec:estimates})}
\begin{itemize}
  \item Why: $\varrho$ is imprecise and close to the lower edge of its
    range. At 0.96 efficiency lasts 25 years.
  \item The horizon roughly doubles (6.2 at 25, 5.4 at 75 at the model's
    moving rate; 6.7, 9.8, 7.0 with the moving rates in UKHLS) and survival
    starts to matter (with loss of efficiency alone, 14.1 at 25 against 7.8
    at 75).
  \item The fit is worse: shares of 43.1, 45.9, 36.8 against 21.0, 16.9,
    9.7 in the data.
  \item The test: do purchases still account for the fall in the inflow?
    \open{Not computed yet: the accounting at 0.96 needs only the existing
    simulation (flows\_by\_age, varrho 0.96), not a new run. The flow,
    horizon, and accounting tables at 0.96 go to the appendix.}
\end{itemize}

\subsection{The window}

\outpara{Where efficient homes come from (Table tab:flows-accounting,
Panel A)}
\begin{itemize}
  \item Per 100 owners a year, by age (25--34, 35--64, 65--75): the stock
    of A--C owners (22.5, 17.7, 12.7) and how it changes.
  \item Inflows by route. Buying an A--C home: 6.9, 4.1, 1.9. Renovating at
    purchase: 0.3, 0.4, 0.3. Renovating in place: 1.2, 1.4, 1.2. Buying is
    the main route at every age and the only one that falls with age.
  \item Because purchases fall, the share of the inflow that comes from
    renovating in place rises with age: 14, 24, and 36 percent. This is the
    claim Section~\ref{sec:identification} makes when it explains how the
    fixed cost shapes the profile.
  \item Outflows. Loss of efficiency is the largest outflow at every age
    (4.6, 3.5, 2.3), then selling to rent and selling to buy a D--G home;
    death matters only at 65--75 (0.24).
  \item Table note: loss of efficiency is 18--21 percent of the stock
    column, not $1-\varrho=16$ percent, because it applies to the stock
    after the year's purchases and sales (checked against the table).
\end{itemize}

\outpara{Which margin produces the fall (Panels B and C)}
\begin{itemize}
  \item Write the inflow as purchases per owner times the A--C share of
    purchases, plus the in-place renovation rate times the share of D--G
    owners. Decompose its fall from 25--34 and from 35--64 to 65--75.
  \item Purchases by renters and D--G owners, per 100 owners, fall from
    15.0 to 10.1 to 4.5 (A--C owners who sell and buy again are left out:
    they are already in the stock). They account for 99 percent of the
    fall from 25--34 and 103 percent from 35--64.
  \item The margins on which the horizon acts are flat, as 5.1 implies:
    the A--C share of purchases is 40--46 percent and the in-place rate
    1.4--1.7 percent of D--G owners at every age. The other margins offset
    one another: together they account for 1 percent of the fall from
    25--34 and $-3$ percent from 35--64 (Panel C). The flat A--C share of
    purchases is the model's version of Fact~3 (recent buyers are about as
    efficient at every age).
  \item With purchases held at the 25--34 rate, the A--C share at 65--75
    would be 25.3 percent instead of 12.7. Holding any other margin moves
    it by at most 1 point.
  \item Caveat: accounting is not causation. Purchases are a choice, and a
    short horizon is itself a reason not to move
    (Section~\ref{sec:data}). The next paragraph tests the window directly.
\end{itemize}

\outpara{Closing the window (Table tab:mechanisms, Panel B)}
\begin{itemize}
  \item With no efficient homes for sale, the gap falls from 9.6 to 1.3
    points and the shares to 6.9, 5.5, 5.6. The window is necessary for
    the gradient. This is the causal counterpart of the accounting.
\end{itemize}

\outpara{Why renovating in place does not replace buying}
\begin{itemize}
  \item The fixed cost is the extra cost of renovating in place in
    Proposition~\ref{prop:window}. Without it, or with it halved, every
    share rises and the profile flattens (ratio 1.36 and 1.57 against
    1.75): owners who no longer buy renovate in place instead. This is the
    role Section~\ref{sec:identification} gives $\xi_f$.
  \item \open{$\xi_f=0$ is the lower edge of the box [0, 0.2]; the halved
    cell is interior. Keep only the halved cell, or say so in the note.}
\end{itemize}

\outpara{What sets how often the old buy: the bequest motive}
\begin{itemize}
  \item A weak bequest motive ($\upsilon=5$): owners at 65--75 buy 23.7
    homes per 100 instead of 6.1, and the gradient reverses (gap $-7.5$).
    A strong one ($\upsilon=200$): they buy 2.8, and the gap widens to 22.3.
  \item The bequest motive matters through purchases, not through the
    value of the heirs' premium. This is the third force of
    Section~\ref{sec:model}, now with a size.
  \item \open{$\upsilon=5$ is the lower edge of the box [5, 500]. An
    interior weak value (e.g.\ 20) needs a new run.}
\end{itemize}

\outpara{Section close (one paragraph)}
\begin{itemize}
  \item Answer to Section~\ref{sec:data}: the gradient is a purchase
    gradient. Fact~3 is what this predicts.
  \item Bridge to Section~\ref{sec:policy}: a policy that lowers the cost
    of renovating in place acts on the margin that 5.2 shows cannot
    replace buying; one that acts at the point of sale acts on the window.
    Section~\ref{sec:policy} compares them.
\end{itemize}

\outpara{Exhibits (planning note, not text)}
\begin{itemize}
  \item Main text, four exhibits: Figure fig:lifecycle-fit (opening);
    Table tab:payoff-horizon (5.1); Table tab:flows-accounting (5.2; Panel
    A flows, Panel B margins, Panel C decomposition; merged 28 October,
    not yet in the paper); Table tab:mechanisms (Panel A comfort and $\varrho$,
    cited in 5.1; Panel B no efficient homes for sale, $\xi_f$, $\upsilon$,
    cited in 5.2; placed in 5.2, where most rows are used).
  \item Appendix: flows, horizon, and accounting at $\varrho=0.96$.
  \item \open{Before building tab:mechanisms: (i) one baseline: the cells
    were run on the GPU (22.4, 17.8, 12.8; gap 9.6), the flows on the CPU
    (22.5, 17.7, 12.7; gap 9.9); re-run the cells on the CPU or report both;
    (ii) ``purchases'' there are all purchases per 100 owners (6.0 at
    65--75), not the 4.5 of the accounting; use one definition; (iii) drop
    the cells at $\varrho=0.99$; the mandate cell goes to
    Section~\ref{sec:policy}.}
\end{itemize}
